\documentclass[11pt,a4paper]{article}
% Shared preamble for RuPhO English problem statements (match T1 2026 style).
% Use: \input{latex/rupho_en_preamble.tex} after \documentclass.
\usepackage[margin=1in]{geometry}
\usepackage{amsmath,amssymb}
\usepackage{graphicx}
\usepackage{float}
\usepackage{enumitem}
\usepackage{microtype}
\usepackage{caption}
\usepackage{tabularx}
\usepackage{hyperref}

\captionsetup{font=small,labelfont=bf,skip=6pt}

\setlist[itemize]{leftmargin=1.2cm}
\setlist[enumerate]{leftmargin=1.2cm}

% One outer box per subproblem: [id | statement | points] (no inner rules)
\newcommand{\problembox}[3]{%
  \par\addvspace{0.42em}%
  \noindent
  \setlength{\fboxsep}{7.5pt}%
  \setlength{\fboxrule}{0.95pt}%
  \fbox{%
    \begin{minipage}{\dimexpr\linewidth-2\fboxsep-2\fboxrule\relax}
    \begin{tabularx}{\linewidth}{@{}>{\raggedright\arraybackslash\bfseries}p{2.1em} @{\hspace{0.9em}} >{\raggedright\arraybackslash}X @{\hspace{0.9em}} >{\raggedleft\arraybackslash\bfseries}p{2.4em}@{}}
    #1 & #2 & #3
    \end{tabularx}
    \end{minipage}%
  }%
  \par\addvspace{0.32em}%
}


\title{\textbf{Thermodynamic black box}}
\author{\normalsize Y21 (2021) --- English translation}
\date{}
\begin{document}
\maketitle

The thermally insulated system consists of two vessels connected by a thin heat-conducting tube. The intrinsic heat capacity of the tube can be neglected; thermal equilibrium in each vessel and along the tube is established much faster than thermal equilibrium between the vessels.

\begin{figure}[H]
    \centering
    \includegraphics[width=0.72\textwidth]{PE3_EN_assets/fig1.jpg}
    \caption{Figure 1.}
\end{figure}

The heat capacity of the body in vessel No. 2 is known: $C_2=1~кJ/{}^\circ\mathrm{C}$. You can set the initial temperature of vessel No. 2 ($T_{20}$, between $0^\circ\mathrm{C}$ and $100^\circ\mathrm{C}$.) and the time $t$ after connecting the vessels with a tube. The program offered to you shows the temperature in vessel No. 2 at time $t$ (set in minutes) after connecting the vessels with a tube ($t=0$).

\problembox{A1}{Determine the initial temperature of vessel No. 1 ($T_{10}$).}{0.40}

\problembox{A2}{Plot a graph of the dependence of the heat capacity of vessel No. 1 on its temperature.}{3.60}

\problembox{A3}{How long will it take for the temperature of vessel No. 1 to increase from $T_{10}$ to $T_{10}+20^\circ\mathrm{C}$ if, instead of vessel No. 2, a tube is connected to a vessel with a constant temperature $100^\circ\mathrm{C}$?}{3.00}

\section*{Program}

exe The executable file for Windows is published on GitHub.

\vspace{1em}
\noindent\textit{Source:} \url{https://pho.rs/p/244}

\end{document}